\Defines

\begin{description}
    \item[Паллета] унифицированная платформа (EURO 1200~$\times$~800~мм), на которой размещаются коробки для складирования и транспортировки.
    \item[Паллетизация] процесс укладки единичных грузовых мест (коробок) на паллету по определённому алгоритму.
    \item[SKU] (Stock Keeping Unit) --- уникальный складской артикул товара, единица учёта на складе.
    \item[Коэффициент заполнения] (\emph{percolation}) --- отношение суммарного объёма размещённых коробок к объёму прямоугольного параллелепипеда, ограниченного паллетой и максимальной достигнутой высотой.
    \item[Минимальный коэффициент опоры] (\emph{min\_support\_ratio}) --- наименьшее среди всех коробок отношение площади основания, опирающейся на нижележащий слой, к полной площади основания.
    \item[Баланс высот] (\emph{height\_balance}) --- метрика равномерности загрузки паллет при мультипаллетной укладке, основанная на коэффициенте вариации высот штабелей.
    \item[Карта высот] (HeightHandler) --- вспомогательная структура данных, хранящая информацию о текущих высотах на поверхности паллеты.
    \item[Геном] (GenomHandler) --- структура кодирования решения, определяющая порядок коробок, их ориентации, предпочтительные позиции и порог опоры.
    \item[Мутация] оператор изменения генома, порождающий новое решение из существующего; в работе реализовано тринадцать операторов.
    \item[JSON] текстовый формат для обмена и хранения данных.
    \item[API] программный интерфейс взаимодействия между компьютерными программами.
    \item[CSV] (\emph{Comma-Separated Values}) --- текстовый формат хранения табличных данных, используемый для входных и выходных файлов.
\end{description}

%%% Local Variables:
%%% mode: latex
%%% TeX-master: "rpz"
%%% End:
